\section{软硬件螺旋上升与数据差距}

前半讲学术史，本章进入产业化路线。王鹤认为，具身智能的软件和硬件是螺旋上升问题：硬件太激进、不成熟，会拖累智能训练和场景落地；软件能力提升后，又会反过来要求更适合采集、执行和交付的硬件。不能把机器人看成“先造一个通用身体，再等大脑学会一切”，也不能只做软件不管本体约束。

访谈描述里有一个关键判断：所有工业机械臂去年的全球总产值只有约 1000 亿 RMB，和理想一家车企产值相当。这个数字的意义是，机器人行业真实商业盘子还不大，不能只靠宏大叙事估值。银河通用的路线是，在已有硬件基础上做相对专用的智能，同时逐步走向越来越通用的智能。这里的“专用”不是放弃通用，而是先找到能产生生产力的场景。

\begin{figure}[H]
\centering
\includegraphics[width=0.96\textwidth]{software-hardware-spiral.png}
\caption{软硬件螺旋上升：硬件过激会拖累智能，智能成熟又反推硬件。自制概念图，依据 01:25:08--01:44:34 对谈内容与视频描述整理。}
\end{figure}

\begin{knowledgebox}{读图：螺旋不是线性流水线}
硬件、本体、数据、模型和场景不是先后完成的模块，而是互相反推。硬件决定能采什么数据、能执行什么动作；模型决定需要什么传感器和动作空间；场景决定成本和可靠性要求。读图时要看循环，而不是把硬件当成一次性采购。
\end{knowledgebox}

\subsection{VLM 为什么弱于 LLM}

本节解释一个核心数据差距。LLM 的文本数据接近“人类在互联网上说过的话”的大规模覆盖；VLM 的视觉数据只是“人眼观测世界”的很小一部分覆盖；VLA 还要额外加入 action 数据，而 action 数据系统采集也就是近两年才开始。换句话说，视觉-语言模型不是因为架构天生低级，而是因为数据覆盖、标注结构和行动反馈都远弱于文本。

这能解释为什么机器人不能简单复用语言模型 scaling law。语言模型训练的是 token 序列，互联网已经提供大量文本；机器人训练的是视觉、语言、动作、状态、失败、恢复和场景之间的耦合，数据天然更少、更贵、更偏。合成数据、遥操作、仿真和出货回流之所以重要，正是因为它们试图补足 action data 的缺口。

\begin{knowledgebox}{术语消化：LLM、VLM、VLA 的数据差别}
\begin{tabularx}{\textwidth}{p{0.18\textwidth}p{0.34\textwidth}X}
\toprule
模型类型 & 数据覆盖 & 关键短板 \\
\midrule
LLM & 互联网文本接近覆盖大量人类表达与知识交换 & 文本不直接包含物理执行成本。 \\
VLM & 图像/视频覆盖只是一小部分人类视觉经验 & 观测缺少完整三维、时间和因果反馈。 \\
VLA & 还需要动作轨迹、关节/末端状态和执行结果 & action data 刚开始系统采集，真实成本极高。 \\
机器人策略 & 需要视觉、语言、动作、失败恢复和场景目标耦合 & 数据必须接地到硬件和真实任务。 \\
\bottomrule
\end{tabularx}
\end{knowledgebox}

讲者强调，不要把目标一步定得过高：一两年内做出完全通用 VLA，从当前数据和产业条件看并不现实。更稳的路径是围绕一个可批量复制的 application 来发展智能，让模型覆盖这个应用内所有必须完成的事情。例如商超和药店场景中的移动、抓取、放置、补货和整理，技能集合有限，但物体、品牌、门店和摆放方式要充分泛化。

\begin{figure}[H]
\centering
\includegraphics[width=0.94\textwidth]{vlm-vs-llm-data.png}
\caption{VLM 为什么弱于 LLM：视觉覆盖与动作数据远少于文本覆盖。自制概念图，依据 01:25:08--01:44:34 对谈内容与视频描述整理。}
\end{figure}

\begin{knowledgebox}{读图：差距在数据覆盖，而不只是模型名字}
左侧 LLM 享受大规模文本覆盖，右侧 VLM/VLA 需要视觉、三维、动作和反馈。VLA 中的 Action Data 是最稀缺的一层：没有动作轨迹和结果，模型很难从“看懂”走向“做成”。
\end{knowledgebox}

\begin{warningbox}{激进硬件方案会拖累智能}
如果硬件不成熟、维护成本高、动作空间不稳定，模型训练会不断适配错误本体，场景交付也会被可靠性吞噬。具身智能不是越通用越好，而是要在当前硬件可承载的范围内找到能产生真实价值的智能。
\end{warningbox}

\subsection{本章小结}

具身智能不是纯软件或纯硬件问题。软件、硬件、数据和场景必须螺旋上升；VLM/VLA 相比 LLM 的差距，核心在于视觉覆盖不足和 action data 稀缺。

